% !TeX root = ../My_thesis.tex
\chapter{Программный код}\label{appendix-MikTeX-TexStudio}							% Заголовок
%\addcontentsline{toc}{chapter}{Second call for chapters to participate in the book Machine learning in analysis of biomedical and socio-economic data}	% Добавляем его в оглавление
pom.xml
\begin{verbatim}
<?xml version="1.0" encoding="UTF-8"?>
<project xmlns="http://maven.apache.org/POM/4.0.0"
         xmlns:xsi="http://www.w3.org/2001/XMLSchema-instance"
         xsi:schemaLocation="http://maven.apache.org/POM/4.0.0 http://maven.apache.org/xsd/maven-4.0.0.xsd">
    <modelVersion>4.0.0</modelVersion>

    <groupId>com.ipfilter</groupId>
    <artifactId>ip-filter-bloom-vs-hashset</artifactId>
    <version>1.0-SNAPSHOT</version>
    <packaging>jar</packaging>

    <properties>
        <maven.compiler.source>17</maven.compiler.source>
        <maven.compiler.target>17</maven.compiler.target>
        <project.build.sourceEncoding>UTF-8</project.build.sourceEncoding>
        
        <junit.version>5.9.3</junit.version>
        <mockito.version>5.3.1</mockito.version>
        <jackson.version>2.15.2</jackson.version>
        <pitest.version>1.16.0</pitest.version>
        <pitest-junit5-plugin.version>1.2.1</pitest-junit5-plugin.version>
    </properties>

    <dependencies>
        <!-- JUnit 5 -->
        <dependency>
            <groupId>org.junit.jupiter</groupId>
            <artifactId>junit-jupiter-api</artifactId>
            <version>${junit.version}</version>
            <scope>test</scope>
        </dependency>
        <dependency>
            <groupId>org.junit.jupiter</groupId>
            <artifactId>junit-jupiter-engine</artifactId>
            <version>${junit.version}</version>
            <scope>test</scope>
        </dependency>
        <dependency>
            <groupId>org.junit.jupiter</groupId>
            <artifactId>junit-jupiter-params</artifactId>
            <version>${junit.version}</version>
            <scope>test</scope>
        </dependency>

        <!-- Hamcrest -->
        <dependency>
            <groupId>org.hamcrest</groupId>
            <artifactId>hamcrest</artifactId>
            <version>2.2</version>
            <scope>test</scope>
        </dependency>

        <!-- Mockito -->
        <dependency>
            <groupId>org.mockito</groupId>
            <artifactId>mockito-core</artifactId>
            <version>${mockito.version}</version>
            <scope>test</scope>
        </dependency>
        <dependency>
            <groupId>org.mockito</groupId>
            <artifactId>mockito-junit-jupiter</artifactId>
            <version>${mockito.version}</version>
            <scope>test</scope>
        </dependency>

        <!-- Jackson (JSON) -->
        <dependency>
            <groupId>com.fasterxml.jackson.core</groupId>
            <artifactId>jackson-databind</artifactId>
            <version>${jackson.version}</version>
        </dependency>
    </dependencies>

    <build>
        <plugins>
            <plugin>
                <groupId>org.apache.maven.plugins</groupId>
                <artifactId>maven-surefire-plugin</artifactId>
                <version>3.1.2</version>
            </plugin>

            <plugin>
                <groupId>org.pitest</groupId>
                <artifactId>pitest-maven</artifactId>
                <version>${pitest.version}</version>
                <dependencies>
                    <dependency>
                        <groupId>org.pitest</groupId>
                        <artifactId>pitest-junit5-plugin</artifactId>
                        <version>${pitest-junit5-plugin.version}</version>
                    </dependency>
                </dependencies>
                <configuration>
                    <targetClasses>
                        <param>com.ipfilter.filter.*</param>
                        <param>com.ipfilter.service.*</param>
                    </targetClasses>
                    <targetTests>
                        <param>com.ipfilter.filter.*Test</param>
                        <param>com.ipfilter.service.*Test</param>
                    </targetTests>
                    <mutators>
                        <mutator>DEFAULTS</mutator>
                    </mutators>
                    <verbose>true</verbose>
                    <outputFormats>
                        <outputFormat>HTML</outputFormat>
                    </outputFormats>
                    <timestampedReports>false</timestampedReports>
                    <failWhenNoMutations>false</failWhenNoMutations>
                    <threads>1</threads>
                </configuration>
            </plugin>

            <plugin>
                <groupId>org.codehaus.mojo</groupId>
                <artifactId>exec-maven-plugin</artifactId>
                <version>3.1.0</version>
                <configuration>
                    <mainClass>com.ipfilter.Main</mainClass>
                </configuration>
            </plugin>
        </plugins>
    </build>
</project>
\end{verbatim}

\noindent \textbf{main.java}
\begin{verbatim}
package com.ipfilter;


import com.ipfilter.model.OutputData;
import com.ipfilter.service.FilterService;

import java.util.Arrays;
import java.util.List;
import java.util.Scanner;

public class Main {
    public static void main(String[] args) {
        FilterService service = new FilterService();
        Scanner scanner = new Scanner(System.in);
        boolean running = true;

        System.out.println("=== IP Blacklist Filter ===");

        while (running) {
            printMenu();
            System.out.print("Выбор: ");
            String choice = scanner.nextLine();

            try {
                switch (choice) {
                    case "1":
                        System.out.print("Путь к JSON файлу (например, input_example1.json): ");
                        String path = scanner.nextLine();
                        service.loadDataFromFile(path);
                        System.out.println("Данные загружены успешно!");
                        break;
                    case "2":
                        System.out.print("Введите IP для проверки: ");
                        String ip = scanner.nextLine();
                        OutputData singleRes = service.processAndCompare(List.of(ip));
                        System.out.println("HashSet: " + singleRes.getResults().get(0).isHashSetResult());
                        System.out.println("Bloom: " + singleRes.getResults().get(0).isBloomFilterResult());
                        break;
                    case "3":
                        // Для примера проверяем фиксированный список, если файл не загружен
                        List<String> testIps = Arrays.asList("192.168.1.1", "8.8.8.8", "10.0.0.1");
                        OutputData res = service.processAndCompare(testIps);
                        System.out.println("Статистика:");
                        System.out.println("Память HS: " + res.getStatistics().getHashSetMemoryBytes() + " b");
                        System.out.println("Память BF: " + res.getStatistics().getBloomFilterMemoryBytes() + " b");
                        System.out.println("Рекомендация: " + res.getRecommendation());

                        System.out.print("Сохранить в output_example.json? (y/n): ");
                        if (scanner.nextLine().equalsIgnoreCase("y")) {
                            service.saveResults(res, "output_example.json");
                            System.out.println("Сохранено.");
                        }
                        break;
                    case "4":
                        System.out.println("Справка: Программа сравнивает точность HashSet и скорость/память Bloom Filter.");
                        break;
                    case "5":
                        running = false;
                        System.out.println("Выход.");
                        break;
                    default:
                        System.out.println("Неверный выбор.");
                }
            } catch (Exception e) {
                System.out.println("Ошибка: " + e.getMessage());
            }
        }
        scanner.close();
    }

    private static void printMenu() {
        System.out.println("\n1. Загрузить blacklist из JSON");
        System.out.println("2. Проверить IP адрес");
        System.out.println("3. Сравнить алгоритмы (Demo)");
        System.out.println("4. Справка");
        System.out.println("5. Выход");
    }
}

\end{verbatim}

IpFilter.java
\begin{verbatim}
package com.ipfilter.filter;

public interface IpFilter {
    void add(String ip);
    boolean contains(String ip);
    long getMemoryUsageBytes();
    void clear();
    int getSize();
}

\end{verbatim}


BloomFilter.java
\begin{verbatim}
package com.ipfilter.filter;

import java.util.BitSet;
import java.util.zip.CRC32;

public class BloomFilter implements IpFilter {
    private final BitSet bitSet;
    private final int size;
    private final int hashFunctionsCount;
    private final long seed;
    private int elementsCount = 0;

    // Конструктор: размер битового массива, кол-во хэш-функций, сид
    public BloomFilter(int size, int hashFunctionsCount, long seed) {
        this.size = size;
        this.hashFunctionsCount = hashFunctionsCount;
        this.seed = seed;
        this.bitSet = new BitSet(size);
    }

    @Override
    public void add(String ip) {
        if (ip == null || ip.isEmpty()) return;
        for (int i = 0; i < hashFunctionsCount; i++) {
            int index = getHash(ip, i) % size;
            bitSet.set(index);
        }
        elementsCount++;
    }

    @Override
    public boolean contains(String ip) {
        if (ip == null) return false;
        for (int i = 0; i < hashFunctionsCount; i++) {
            int index = getHash(ip, i) % size;
            if (!bitSet.get(index)) {
                return false; // Точно нет
            }
        }
        return true; // Возможно есть (False Positive possible)
    }

    // Детерминированная хэш-функция на основе CRC32 и сида
    private int getHash(String ip, int iteration) {
        CRC32 crc = new CRC32();
        String input = ip + "_" + iteration + "_" + seed;
        crc.update(input.getBytes());
        return Math.abs((int) crc.getValue()); // Всегда положительное!
    }

    @Override
    public long getMemoryUsageBytes() {
        return size / 8L + 32L; // Биты в байты + оверхед объекта
    }

    @Override
    public void clear() {
        bitSet.clear();
        elementsCount = 0;
    }

    @Override
    public int getSize() {
        return elementsCount;
    }

    // Расчет вероятности False Positive: (1 - e^(-kn/m))^k
    public double getFalsePositiveProbability() {
        if (size == 0 || elementsCount == 0) return 0.0;
        double m = size;
        double k = hashFunctionsCount;
        double n = elementsCount;
        return Math.pow((1 - Math.exp(-k * n / m)), k);
    }
}

\end{verbatim}

HashSetFilter.java
\begin{verbatim}
package com.ipfilter.filter;

import java.util.HashSet;
import java.util.Set;

public class HashSetFilter implements IpFilter {
    private final Set<String> blacklist;

    public HashSetFilter() {
        this.blacklist = new HashSet<>();
    }

    @Override
    public void add(String ip) {
        if (ip != null && !ip.isEmpty()) {
            blacklist.add(ip);
        }
    }

    @Override
    public boolean contains(String ip) {
        // Branch Testing point: null check
        if (ip == null) {
            return false;
        }
        return blacklist.contains(ip);
    }

    @Override
    public long getMemoryUsageBytes() {
        // Approximation: Object overhead + String size + HashSet overhead
        // Roughly 64 bytes per entry + string content
        return blacklist.size() * 100L;
    }

    @Override
    public void clear() {
        blacklist.clear();
    }

    @Override
    public int getSize() {
        return blacklist.size();
    }
}

\end{verbatim}

FilterResult.java
\begin{verbatim}
package com.ipfilter.model;

import com.fasterxml.jackson.annotation.JsonProperty;

public class FilterResult {
    @JsonProperty("ip")
    private String ip;
    @JsonProperty("hashSetResult")
    private boolean hashSetResult;
    @JsonProperty("bloomFilterResult")
    private boolean bloomFilterResult;
    @JsonProperty("match")
    private boolean match; // Совпадают ли результаты

    public FilterResult() {}

    public FilterResult(String ip, boolean hashSetResult, boolean bloomFilterResult) {
        this.ip = ip;
        this.hashSetResult = hashSetResult;
        this.bloomFilterResult = bloomFilterResult;
        this.match = (hashSetResult == bloomFilterResult);
    }

    // Getters and Setters
    public String getIp() { return ip; }
    public void setIp(String ip) { this.ip = ip; }
    public boolean isHashSetResult() { return hashSetResult; }
    public void setHashSetResult(boolean hashSetResult) { this.hashSetResult = hashSetResult; }
    public boolean isBloomFilterResult() { return bloomFilterResult; }
    public void setBloomFilterResult(boolean bloomFilterResult) { this.bloomFilterResult = bloomFilterResult; }
    public boolean isMatch() { return match; }
    public void setMatch(boolean match) { this.match = match; }
}
\end{verbatim}

FilterStatistics.java
\begin{verbatim}
package com.ipfilter.model;

import com.fasterxml.jackson.annotation.JsonProperty;

public class FilterStatistics {
    @JsonProperty("hashSetMemoryBytes")
    private long hashSetMemoryBytes;
    @JsonProperty("bloomFilterMemoryBytes")
    private long bloomFilterMemoryBytes;
    @JsonProperty("hashSetTimeMs")
    private long hashSetTimeMs;
    @JsonProperty("bloomFilterTimeMs")
    private long bloomFilterTimeMs;
    @JsonProperty("falsePositiveProbability")
    private double falsePositiveProbability;

    // Getters and Setters
    public long getHashSetMemoryBytes() { return hashSetMemoryBytes; }
    public void setHashSetMemoryBytes(long hashSetMemoryBytes) { this.hashSetMemoryBytes = hashSetMemoryBytes; }
    public long getBloomFilterMemoryBytes() { return bloomFilterMemoryBytes; }
    public void setBloomFilterMemoryBytes(long bloomFilterMemoryBytes) { this.bloomFilterMemoryBytes = bloomFilterMemoryBytes; }
    public long getHashSetTimeMs() { return hashSetTimeMs; }
    public void setHashSetTimeMs(long hashSetTimeMs) { this.hashSetTimeMs = hashSetTimeMs; }
    public long getBloomFilterTimeMs() { return bloomFilterTimeMs; }
    public void setBloomFilterTimeMs(long bloomFilterTimeMs) { this.bloomFilterTimeMs = bloomFilterTimeMs; }
    public double getFalsePositiveProbability() { return falsePositiveProbability; }
    public void setFalsePositiveProbability(double falsePositiveProbability) { this.falsePositiveProbability = falsePositiveProbability; }
}
\end{verbatim}

InputData.java
\begin{verbatim}
package com.ipfilter.model;

import com.fasterxml.jackson.annotation.JsonProperty;
import java.util.List;

public class InputData {
    @JsonProperty("blacklist")
    private List<String> blacklist;
    @JsonProperty("checkIps")
    private List<String> checkIps;

    public List<String> getBlacklist() { return blacklist; }
    public void setBlacklist(List<String> blacklist) { this.blacklist = blacklist; }
    public List<String> getCheckIps() { return checkIps; }
    public void setCheckIps(List<String> checkIps) { this.checkIps = checkIps; }
}
\end{verbatim}

OutputData.java
\begin{verbatim}
package com.ipfilter.model;

import com.fasterxml.jackson.annotation.JsonProperty;
import java.util.List;

public class OutputData {
    @JsonProperty("timestamp")
    private String timestamp;
    @JsonProperty("results")
    private List<FilterResult> results;
    @JsonProperty("statistics")
    private FilterStatistics statistics;
    @JsonProperty("recommendation")
    private String recommendation;

    // Getters and Setters
    public String getTimestamp() { return timestamp; }
    public void setTimestamp(String timestamp) { this.timestamp = timestamp; }
    public List<FilterResult> getResults() { return results; }
    public void setResults(List<FilterResult> results) { this.results = results; }
    public FilterStatistics getStatistics() { return statistics; }
    public void setStatistics(FilterStatistics statistics) { this.statistics = statistics; }
    public String getRecommendation() { return recommendation; }
    public void setRecommendation(String recommendation) { this.recommendation = recommendation; }
}
\end{verbatim}

FilterService.java
\begin{verbatim}
package com.ipfilter.model;

import com.fasterxml.jackson.annotation.JsonProperty;
import java.util.List;

public class OutputData {
    @JsonProperty("timestamp")
    private String timestamp;
    @JsonProperty("results")
    private List<FilterResult> results;
    @JsonProperty("statistics")
    private FilterStatistics statistics;
    @JsonProperty("recommendation")
    private String recommendation;

    // Getters and Setters
    public String getTimestamp() { return timestamp; }
    public void setTimestamp(String timestamp) { this.timestamp = timestamp; }
    public List<FilterResult> getResults() { return results; }
    public void setResults(List<FilterResult> results) { this.results = results; }
    public FilterStatistics getStatistics() { return statistics; }
    public void setStatistics(FilterStatistics statistics) { this.statistics = statistics; }
    public String getRecommendation() { return recommendation; }
    public void setRecommendation(String recommendation) { this.recommendation = recommendation; }
}
\end{verbatim}

JsonHandler.java
\begin{verbatim}
package com.ipfilter.service;

import com.fasterxml.jackson.databind.ObjectMapper;
import com.ipfilter.model.InputData;
import com.ipfilter.model.OutputData;

import java.io.File;
import java.io.IOException;
import java.io.InputStream;

public class JsonHandler {
    private final ObjectMapper objectMapper;

    public JsonHandler() {
        this.objectMapper = new ObjectMapper();
        this.objectMapper.findAndRegisterModules();
    }

    // Файл: src/main/java/com/ipfilter/service/JsonHandler.java

    public InputData loadInput(String filePath) throws IOException {
        // Сначала пробуем загрузить из classpath (для тестов)
        InputStream inputStream = getClass().getClassLoader().getResourceAsStream(filePath);

        if (inputStream != null) {
            return objectMapper.readValue(inputStream, InputData.class);
        }

        // Если не нашли в classpath, пробуем как обычный файл (для production)
        File file = new File(filePath);
        if (file.exists()) {
            return objectMapper.readValue(file, InputData.class);
        }

        throw new IOException("File not found: " + filePath);
    }

    public void saveOutput(String filePath, OutputData data) throws IOException {
        objectMapper.writerWithDefaultPrettyPrinter()
                .writeValue(new File(filePath), data);
    }
}
\end{verbatim}

BloomFilterTest.java
\begin{verbatim}
package com.ipfilter.filter;

import org.junit.jupiter.api.BeforeEach;
import org.junit.jupiter.api.Test;
import org.junit.jupiter.api.Assumptions;
import org.junit.jupiter.params.ParameterizedTest;
import org.junit.jupiter.params.provider.CsvSource;

import static org.junit.jupiter.api.Assertions.*;
import static org.hamcrest.MatcherAssert.assertThat;
import static org.hamcrest.Matchers.*;

class BloomFilterTest {

    private BloomFilter filter;

    @BeforeEach
    void setUp() {
        // Seed 42 для детерминированности
        filter = new BloomFilter(1000, 3, 42L);
    }

    // --- SPECIFICATION BASED TESTING ---

    // Метод 1: Boundary Value Analysis (BVA)
    @Test
    void testBoundary_FullCapacityBehavior() {
        // Заполняем фильтр (меньше чем размер битового массива)
        for (int i = 0; i < 100; i++) {
            filter.add("IP_" + i);
        }
        double fpProb = filter.getFalsePositiveProbability();
        assertThat("Вероятность FP должна быть >= 0", fpProb, greaterThanOrEqualTo(0.0));
    }

    // Метод 2: Equivalence Partitioning (EP)
    @ParameterizedTest
    @CsvSource({
            "192.168.1.1, true",
            "10.0.0.1, true"
    })
    void testEquivalenceClasses(String ip, boolean shouldBeAdded) {
        filter.add(ip);
        // Bloom Filter не даёт False Negative
        assertTrue(filter.contains(ip), "Добавленный элемент должен быть найден");
    }

    // --- STRUCTURE BASED TESTING ---

    // Метод 1: Branch Testing
    @Test
    void testBranch_NotFoundReturnsFalse() {
        filter.add("1.1.1.1");
        // Элемент который точно не добавлен
        boolean result = filter.contains("2.2.2.2");
        assertNotNull(result);
    }

    // Метод 2: Data Flow Testing
    @Test
    void testDataFlow_ElementsCountTracking() {
        assertEquals(0, filter.getSize());
        filter.add("A");
        assertEquals(1, filter.getSize());
        filter.add("B");
        assertEquals(2, filter.getSize());
        filter.clear();
        assertEquals(0, filter.getSize());
    }

    // --- JUNIT REQUIREMENTS (Assumptions & Matchers) ---

    @Test
    void testAssumptions_Environment() {
        Assumptions.assumeTrue(filter.getMemoryUsageBytes() > 0, "Память должна быть выделена");

        filter.add("TestIP");

        Assumptions.assumingThat(filter.contains("TestIP"), () -> {
            double prob = filter.getFalsePositiveProbability();
            assertThat(prob, greaterThanOrEqualTo(0.0));
        });
    }

    @Test
    void testHamcrestMatchers() {
        filter.add("MatchTest");
        assertThat(filter.contains("MatchTest"), equalTo(true));
        assertThat(filter.getFalsePositiveProbability(), lessThan(1.0));
    }

    // === НОВЫЕ ТЕСТЫ ДЛЯ MUTATION COVERAGE ===

    @Test
    void testAdd_NullValue() {
        filter.add(null);
        assertEquals(0, filter.getSize());
    }

    @Test
    void testAdd_EmptyString() {
        filter.add("");
        assertEquals(0, filter.getSize());
    }

    @Test
    void testContains_NullInput() {
        assertFalse(filter.contains(null));
    }

    @Test
    void testClear_ResetsCount() {
        filter.add("1.1.1.1");
        filter.add("2.2.2.2");
        filter.clear();
        assertEquals(0, filter.getSize());
        assertEquals(0.0, filter.getFalsePositiveProbability());
    }

    @Test
    void testFalsePositiveProbability_EmptyFilter() {
        assertEquals(0.0, filter.getFalsePositiveProbability());
    }

    @Test
    void testFalsePositiveProbability_GrowsWithElements() {
        double probBefore = filter.getFalsePositiveProbability();
        for (int i = 0; i < 50; i++) {
            filter.add("IP_" + i);
        }
        double probAfter = filter.getFalsePositiveProbability();
        assertTrue(probAfter > probBefore, "Вероятность FP должна расти");
    }

    @Test
    void testMemoryUsage_Calculation() {
        long memory = filter.getMemoryUsageBytes();
        assertTrue(memory > 0, "Память должна быть > 0");
    }

    @Test
    void testNoFalseNegatives() {
        // Bloom Filter НЕ должен давать False Negative
        String testIp = "192.168.1.1";
        filter.add(testIp);
        assertTrue(filter.contains(testIp), "False Negative недопустим!");
    }
    @Test
    void testAdd_Null_DoesNotIncreaseSize() {
        filter.add(null);
        assertEquals(0, filter.getSize());
    }

    @Test
    void testAdd_EmptyString_DoesNotIncreaseSize() {
        filter.add("");
        assertEquals(0, filter.getSize());
    }

    @Test
    void testContains_Null_ReturnsFalse() {
        assertFalse(filter.contains(null));
    }

    @Test
    void testClear_ResetsEverything() {
        filter.add("1.1.1.1");
        filter.add("2.2.2.2");
        assertEquals(2, filter.getSize());
        filter.clear();
        assertEquals(0, filter.getSize());
        assertEquals(0.0, filter.getFalsePositiveProbability(), 0.001);
    }

    @Test
    void testFalsePositiveProbability_IncreasesWithElements() {
        double probBefore = filter.getFalsePositiveProbability();
        for (int i = 0; i < 100; i++) {
            filter.add("IP_" + i);
        }
        double probAfter = filter.getFalsePositiveProbability();
        assertTrue(probAfter > probBefore, "Вероятность FP должна расти с добавлением элементов");
    }

    @Test
    void testMemoryUsage_AlwaysPositive() {
        long memory = filter.getMemoryUsageBytes();
        assertTrue(memory > 0, "Память должна быть положительной");
    }

    @Test
    void testNoFalseNegatives_Ever() {
        String[] testIps = {"1.1.1.1", "2.2.2.2", "3.3.3.3", "192.168.1.1", "10.0.0.1"};
        for (String ip : testIps) {
            filter.add(ip);
            assertTrue(filter.contains(ip), "False Negative недопустим для IP: " + ip);
        }
    }

    @Test
    void testMultipleHashFunctions_WorkCorrectly() {
        filter.add("test.ip");
        // Bloom Filter должен найти добавленный элемент
        assertTrue(filter.contains("test.ip"));
    }

    @Test
    void testHashFunction_DifferentIterations() {
        // Тестируем что разные итерации дают разные хэши
        filter.add("test.ip");
        assertTrue(filter.contains("test.ip"));
    }

    @Test
    void testBitSetOperations() {
        // Тестируем установку битов
        BloomFilter smallFilter = new BloomFilter(100, 3, 42L);
        smallFilter.add("1.1.1.1");
        assertTrue(smallFilter.contains("1.1.1.1"));
    }

    @Test
    void testFalsePositiveProbability_ZeroElements() {
        assertEquals(0.0, filter.getFalsePositiveProbability(), 0.001);
    }

    @Test
    void testFalsePositiveProbability_AfterClear() {
        filter.add("1.1.1.1");
        filter.add("2.2.2.2");
        filter.clear();
        assertEquals(0.0, filter.getFalsePositiveProbability(), 0.001);
    }

    @Test
    void testSize_AfterClear() {
        filter.add("1.1.1.1");
        assertEquals(1, filter.getSize());
        filter.clear();
        assertEquals(0, filter.getSize());
    }

    @Test
    void testContains_NotAddedIP() {
        // IP который точно не добавлен
        assertFalse(filter.contains("99.99.99.99"));
    }

    @Test
    void testAdd_MultipleSameIP() {
        filter.add("1.1.1.1");
        filter.add("1.1.1.1");
        filter.add("1.1.1.1");
        // Bloom Filter не отслеживает дубликаты, но должен найти
        assertTrue(filter.contains("1.1.1.1"));
    }

    @Test
    void testDifferentSeeds_ProduceDifferentResults() {
        BloomFilter filter1 = new BloomFilter(1000, 3, 42L);
        BloomFilter filter2 = new BloomFilter(1000, 3, 100L);

        filter1.add("test.ip");
        filter2.add("test.ip");

        // Оба должны найти добавленный IP
        assertTrue(filter1.contains("test.ip"));
        assertTrue(filter2.contains("test.ip"));
    }

    @Test
    void testLargeNumberOfElements() {
        // Добавляем много элементов
        for (int i = 0; i < 200; i++) {
            filter.add("IP_" + i);
        }
        assertEquals(200, filter.getSize());

        // Проверяем что все находятся (no false negatives)
        for (int i = 0; i < 200; i++) {
            assertTrue(filter.contains("IP_" + i));
        }
    }

    @Test
    void testProbability_GrowsWithElements() {
        double probBefore = filter.getFalsePositiveProbability();

        for (int i = 0; i < 100; i++) {
            filter.add("IP_" + i);
        }

        double probAfter = filter.getFalsePositiveProbability();
        assertTrue(probAfter > probBefore, "Вероятность FP должна расти");
    }
}
\end{verbatim}

HashSetFilterTest.java
\begin{verbatim}
package com.ipfilter.filter;

import org.junit.jupiter.api.BeforeEach;
import org.junit.jupiter.api.Test;
import org.junit.jupiter.params.ParameterizedTest;
import org.junit.jupiter.params.provider.ValueSource;

import static org.junit.jupiter.api.Assertions.*;
import static org.hamcrest.MatcherAssert.assertThat;

class HashSetFilterTest {

    private IpFilter filter;

    @BeforeEach
    void setUp() {
        filter = new HashSetFilter();
    }

    // --- SPECIFICATION BASED TESTING ---

    // Метод 1: Boundary Value Analysis (BVA)
    @Test
    void testBoundaryValues_EmptyAndSingleElement() {
        // Граница: Пустой фильтр
        assertFalse(filter.contains("192.168.1.1"), "Пустой фильтр не должен содержать IP");

        // Граница: Добавление 1 элемента
        filter.add("192.168.1.1");
        assertTrue(filter.contains("192.168.1.1"), "Фильтр должен содержать добавленный IP");

        // Граница: Проверка несуществующего рядом с существующим
        assertFalse(filter.contains("192.168.1.2"), "Фильтр не должен содержать похожий IP");
    }

    // Метод 2: Equivalence Partitioning (EP)
    @ParameterizedTest
    @ValueSource(strings = {"10.0.0.1", "172.16.0.1", "8.8.8.8"})
    void testEquivalencePartitioning_ValidIps(String ip) {
        // Класс эквивалентности: Валидные IP
        filter.add(ip);
        assertTrue(filter.contains(ip));
    }

    @Test
    void testEquivalencePartitioning_InvalidInputs() {
        // Класс эквивалентности: Null и пустые строки (не должны ломать код)
        filter.add(null);
        filter.add("");
        assertEquals(0, filter.getSize(), "Размер не должен измениться при добавлении невалидных данных");
    }

    // --- STRUCTURE BASED TESTING ---

    // Метод 1: Branch Testing
    @Test
    void testBranchCoverage_NullInputContains() {
        // Проверяем ветку if (ip == null) внутри contains
        boolean result = filter.contains(null);
        assertFalse(result, "Ветка обработки null должна возвращать false");
    }

    // Метод 2: Data Flow Testing
    @Test
    void testDataFlow_DefinitionUse() {
        // Def: создание фильтра
        IpFilter localFilter = new HashSetFilter();
        // Use: добавление
        localFilter.add("1.1.1.1");
        // Use: проверка и получение размера
        int size = localFilter.getSize();
        boolean exists = localFilter.contains("1.1.1.1");

        assertTrue(exists);
        assertEquals(1, size);
    }

    // --- JUNIT REQUIREMENTS ---

    @Test
    void testAssertionsVariety() {
        filter.add("1.1.1.1");
        filter.add("2.2.2.2");

        // 1. assertEquals
        assertEquals(2, filter.getSize());
        // 2. assertTrue
        assertTrue(filter.contains("1.1.1.1"));
        // 3. assertFalse
        assertFalse(filter.contains("3.3.3.3"));
        // 4. assertThrows (проверка на отсутствие исключений при корректной работе, либо тестирование edge cases)
        assertDoesNotThrow(() -> filter.clear());
        // 5. assertAll (группировка)
        assertAll("Filter State",
                () -> assertEquals(0, filter.getSize()),
                () -> assertFalse(filter.contains("1.1.1.1"))
        );
    }

    // === НОВЫЕ ТЕСТЫ ДЛЯ MUTATION COVERAGE ===

    @Test
    void testAdd_NullValue() {
        // Тест на обработку null
        filter.add(null);
        assertEquals(0, filter.getSize(), "Null не должен добавляться");
    }

    @Test
    void testAdd_EmptyString() {
        // Тест на обработку пустой строки
        filter.add("");
        assertEquals(0, filter.getSize(), "Пустая строка не должна добавляться");
    }

    @Test
    void testAdd_DuplicateValues() {
        // Тест на дубликаты
        filter.add("192.168.1.1");
        filter.add("192.168.1.1");
        assertEquals(1, filter.getSize(), "Дубликаты не должны увеличивать размер");
    }

    @Test
    void testContains_NullInput() {
        // Тест ветки if (ip == null)
        assertFalse(filter.contains(null));
    }

    @Test
    void testClear() {
        // Тест метода clear
        filter.add("1.1.1.1");
        filter.add("2.2.2.2");
        filter.clear();
        assertEquals(0, filter.getSize());
        assertFalse(filter.contains("1.1.1.1"));
    }

    @Test
    void testMemoryUsage() {
        // Тест расчёта памяти
        filter.add("1.1.1.1");
        long memoryBefore = filter.getMemoryUsageBytes();
        filter.add("2.2.2.2");
        long memoryAfter = filter.getMemoryUsageBytes();
        assertTrue(memoryAfter > memoryBefore, "Память должна расти");
    }

    @Test
    void testAdd_Null_DoesNotIncreaseSize() {
        filter.add(null);
        assertEquals(0, filter.getSize(), "Null не должен увеличивать размер");
    }

    @Test
    void testAdd_EmptyString_DoesNotIncreaseSize() {
        filter.add("");
        assertEquals(0, filter.getSize(), "Пустая строка не должна увеличивать размер");
    }

    @Test
    void testAdd_Duplicate_DoesNotIncreaseSize() {
        filter.add("192.168.1.1");
        filter.add("192.168.1.1");
        assertEquals(1, filter.getSize(), "Дубликаты не должны увеличивать размер");
    }

    @Test
    void testContains_Null_ReturnsFalse() {
        assertFalse(filter.contains(null), "Null должен возвращать false");
    }

    @Test
    void testClear_ResetsSize() {
        filter.add("1.1.1.1");
        filter.add("2.2.2.2");
        assertEquals(2, filter.getSize());
        filter.clear();
        assertEquals(0, filter.getSize(), "Clear должен обнулить размер");
        assertFalse(filter.contains("1.1.1.1"), "После clear элементы не должны находиться");
    }

    @Test
    void testMemoryUsage_GrowsWithSize() {
        long memoryBefore = filter.getMemoryUsageBytes();
        filter.add("1.1.1.1");
        long memoryAfter = filter.getMemoryUsageBytes();
        assertTrue(memoryAfter >= memoryBefore, "Память должна расти или оставаться той же");
    }

    @Test
    void testMultipleAdds_DifferentIps() {
        filter.add("1.1.1.1");
        filter.add("2.2.2.2");
        filter.add("3.3.3.3");
        assertEquals(3, filter.getSize());
        assertTrue(filter.contains("1.1.1.1"));
        assertTrue(filter.contains("2.2.2.2"));
        assertTrue(filter.contains("3.3.3.3"));
        assertFalse(filter.contains("4.4.4.4"));
    }

    @Test
    void testMemoryUsage_EmptyFilter() {
        long memory = filter.getMemoryUsageBytes();
        assertEquals(0, memory, "Память пустого фильтра должна быть 0");
    }

    @Test
    void testMemoryUsage_AfterAdd() {
        filter.add("1.1.1.1");
        long memory = filter.getMemoryUsageBytes();
        assertTrue(memory > 0, "Память должна быть > 0 после добавления");
    }

    @Test
    void testSize_AfterMultipleAdds() {
        filter.add("1.1.1.1");
        filter.add("2.2.2.2");
        filter.add("3.3.3.3");
        assertEquals(3, filter.getSize());
    }

    @Test
    void testSize_AfterClear() {
        filter.add("1.1.1.1");
        filter.clear();
        assertEquals(0, filter.getSize());
    }

    @Test
    void testContains_AfterClear() {
        filter.add("1.1.1.1");
        filter.clear();
        assertFalse(filter.contains("1.1.1.1"));
    }

    @Test
    void testAdd_NullDoesNotIncreaseSize() {
        int sizeBefore = filter.getSize();
        filter.add(null);
        assertEquals(sizeBefore, filter.getSize());
    }

    @Test
    void testAdd_EmptyStringDoesNotIncreaseSize() {
        int sizeBefore = filter.getSize();
        filter.add("");
        assertEquals(sizeBefore, filter.getSize());
    }

    @Test
    void testContains_NullReturnsFalse() {
        assertFalse(filter.contains(null));
    }

    @Test
    void testClear_MultipleTimes() {
        filter.add("1.1.1.1");
        filter.clear();
        filter.clear(); // Второй вызов не должен ломать
        assertEquals(0, filter.getSize());
    }

    @Test
    void testAdd_SpecialCharacters() {
        filter.add("test@ip#format");
        assertTrue(filter.contains("test@ip#format"));
    }

    @Test
    void testAdd_VeryLongIP() {
        String longIp = "192.168.1.1";
        filter.add(longIp);
        assertTrue(filter.contains(longIp));
    }

    @Test
    void testContains_CaseSensitive() {
        filter.add("192.168.1.1");
        assertFalse(filter.contains("192.168.1.2"));
    }
}
\end{verbatim}

FilterServiceTest.java
\begin{verbatim}
package com.ipfilter.service;

import com.ipfilter.filter.IpFilter;
import com.ipfilter.model.FilterResult;
import com.ipfilter.model.FilterStatistics;
import com.ipfilter.model.OutputData;
import org.junit.jupiter.api.BeforeEach;
import org.junit.jupiter.api.Test;
import org.junit.jupiter.api.Assumptions;
import org.junit.jupiter.params.ParameterizedTest;
import org.junit.jupiter.params.provider.ValueSource;
import org.junit.jupiter.params.provider.CsvSource;

import java.io.IOException;
import java.util.ArrayList;
import java.util.List;

import static org.junit.jupiter.api.Assertions.*;
import static org.hamcrest.MatcherAssert.assertThat;
import static org.hamcrest.Matchers.*;
import static org.mockito.Mockito.*;

/**
 * Тесты для FilterService с максимальным покрытием мутаций
 *
 * Используемые техники тестирования:
 * - Specification-Based: Boundary Value Analysis, Equivalence Partitioning
 * - Structure-Based: Branch Testing, Data Flow Testing
 */
class FilterServiceTest {

    private FilterService service;

    @BeforeEach
    void setUp() {
        service = new FilterService();
    }

    // ========================================================================
    // SPECIFICATION-BASED TESTING
    // ========================================================================

    /**
     * Метод 1: Boundary Value Analysis (BVA)
     * Тестируем граничные значения количества IP
     */
    @Test
    void testBoundary_EmptyList() {
        // Граница: 0 элементов
        OutputData result = service.processAndCompare(new ArrayList<>());

        assertNotNull(result);
        assertEquals(0, result.getResults().size(), "Пустой список должен вернуть 0 результатов");
        assertNotNull(result.getTimestamp());
        assertNotNull(result.getStatistics());
        assertNotNull(result.getRecommendation());
    }

    @Test
    void testBoundary_SingleElement() {
        // Граница: 1 элемент
        OutputData result = service.processAndCompare(List.of("192.168.1.1"));

        assertEquals(1, result.getResults().size(), "Должен быть 1 результат");
        assertNotNull(result.getResults().get(0).getIp());
    }

    @Test
    void testBoundary_ManyElements() {
        // Граница: много элементов
        List<String> ips = new ArrayList<>();
        for (int i = 0; i < 100; i++) {
            ips.add(i + ".0.0." + i);
        }

        OutputData result = service.processAndCompare(ips);

        assertEquals(100, result.getResults().size());
        assertNotNull(result.getStatistics());
    }

    /**
     * Метод 2: Equivalence Partitioning (EP)
     * Тестируем классы эквивалентности IP адресов
     */
    @ParameterizedTest
    @ValueSource(strings = {
            "192.168.1.1",
            "10.0.0.1",
            "172.16.0.1",
            "8.8.8.8",
            "1.1.1.1"
    })
    void testEquivalence_ValidIpAddresses(String ip) {
        OutputData result = service.processAndCompare(List.of(ip));

        assertNotNull(result);
        assertEquals(1, result.getResults().size());
        assertEquals(ip, result.getResults().get(0).getIp());
    }

    @ParameterizedTest
    @CsvSource({
            "invalid.ip",
            "256.256.256.256",
            "abc.def.ghi.jkl"
    })
    void testEquivalence_InvalidIpAddresses(String ip) {
        assertDoesNotThrow(() -> {
            OutputData result = service.processAndCompare(List.of(ip));
            assertNotNull(result);
        });
    }

    // ========================================================================
    // STRUCTURE-BASED TESTING
    // ========================================================================

    /**
     * Метод 1: Branch Testing
     * Покрываем все ветви if/else в коде
     */
    @Test
    void testBranch_ResultSetMatchTrue() {
        // Ветка: оба фильтра возвращают false → match = true
        OutputData result = service.processAndCompare(List.of("8.8.8.8"));

        FilterResult filterResult = result.getResults().get(0);
        assertFalse(filterResult.isHashSetResult());
        assertFalse(filterResult.isBloomFilterResult());
        assertTrue(filterResult.isMatch(), "Результаты должны совпадать");
    }

    @Test
    void testBranch_StatisticsFieldsPopulated() {
        // Ветка: заполнение статистики
        OutputData result = service.processAndCompare(List.of("1.1.1.1"));
        FilterStatistics stats = result.getStatistics();

        // Проверяем что все поля заполнены и валидны
        assertTrue(stats.getHashSetMemoryBytes() >= 0, "HashSet память должна быть >= 0");
        assertTrue(stats.getBloomFilterMemoryBytes() >= 0, "Bloom память должна быть >= 0");
        assertTrue(stats.getHashSetTimeMs() >= 0, "HashSet время должно быть >= 0");
        assertTrue(stats.getBloomFilterTimeMs() >= 0, "Bloom время должно быть >= 0");
        assertTrue(stats.getFalsePositiveProbability() >= 0, "Вероятность FP должна быть >= 0");
        assertTrue(stats.getFalsePositiveProbability() <= 1, "Вероятность FP должна быть <= 1");
    }

    @Test
    void testBranch_RecommendationGenerated() {
        // Ветка: генерация рекомендации
        OutputData result = service.processAndCompare(List.of("1.1.1.1"));

        assertNotNull(result.getRecommendation());
        assertFalse(result.getRecommendation().isEmpty());
        assertThat(result.getRecommendation().length(), greaterThan(10));
    }

    /**
     * Метод 2: Data Flow Testing
     * Отслеживаем определение и использование переменных
     */
    @Test
    void testDataFlow_OutputDataCreation() {
        // Def: создание OutputData
        OutputData result = service.processAndCompare(List.of("192.168.1.1"));

        // Use: проверка всех полей
        assertNotNull(result.getTimestamp(), "Timestamp должен быть определён");
        assertNotNull(result.getResults(), "Results должны быть определены");
        assertNotNull(result.getStatistics(), "Statistics должны быть определены");
        assertNotNull(result.getRecommendation(), "Recommendation должен быть определён");

        // Use: проверка вложенных объектов
        assertEquals(1, result.getResults().size());
        assertNotNull(result.getResults().get(0).getIp());
    }

    @Test
    void testDataFlow_FilterResultCreation() {
        // Def: создание FilterResult
        OutputData output = service.processAndCompare(List.of("10.0.0.1"));
        FilterResult filterResult = output.getResults().get(0);

        // Use: проверка полей
        assertEquals("10.0.0.1", filterResult.getIp());
        // match вычисляется из hashSetResult и bloomFilterResult
        assertEquals(
                filterResult.isHashSetResult() == filterResult.isBloomFilterResult(),
                filterResult.isMatch()
        );
    }

    // ========================================================================
    // JUNIT 5 REQUIREMENTS
    // ========================================================================

    /**
     * Минимум 5 различных методов assert
     */
    @Test
    void testAssertionsVariety() {
        OutputData result = service.processAndCompare(List.of("1.1.1.1", "2.2.2.2"));

        // 1. assertEquals
        assertEquals(2, result.getResults().size());

        // 2. assertTrue
        assertTrue(result.getResults().size() > 0);

        // 3. assertFalse
        assertFalse(result.getResults().isEmpty());

        // 4. assertNotNull
        assertNotNull(result.getTimestamp());

        // 5. assertAll (группировка ассертов)
        assertAll("OutputData validation",
                () -> assertNotNull(result.getResults()),
                () -> assertNotNull(result.getStatistics()),
                () -> assertNotNull(result.getRecommendation()),
                () -> assertEquals(2, result.getResults().size())
        );

        // 6. assertDoesNotThrow
        assertDoesNotThrow(() -> service.processAndCompare(List.of("3.3.3.3")));
    }

    /**
     * Минимум 2 различных метода assume
     */
    @Test
    void testAssumptionsUsage() {
        // Assume 1: assumeTrue
        Assumptions.assumeTrue(service != null, "Сервис должен быть создан");

        OutputData result = service.processAndCompare(List.of("1.1.1.1"));

        // Assume 2: assumingThat
        Assumptions.assumingThat(
                result.getResults().size() > 0,
                () -> {
                    assertNotNull(result.getResults().get(0).getIp());
                }
        );

        // Assume 3: assumeFalse
        Assumptions.assumeFalse(result.getResults().isEmpty(), "Результаты не должны быть пустыми");
    }

    /**
     * Parameterized Tests
     */
    @ParameterizedTest
    @ValueSource(ints = {1, 5, 10, 50})
    void testParameterized_DifferentIpCounts(int count) {
        List<String> ips = new ArrayList<>();
        for (int i = 0; i < count; i++) {
            ips.add(i + ".0.0." + i);
        }

        OutputData result = service.processAndCompare(ips);

        assertEquals(count, result.getResults().size());
        assertNotNull(result.getStatistics());
        assertNotNull(result.getRecommendation());
    }

    // ========================================================================
    // MOCKITO REQUIREMENTS
    // ========================================================================

    /**
     * Mock Type 1: Standard Mock behavior (when/thenReturn)
     */
    @Test
    void testMockBehavior_VerifyWithLenient() {
        // Создаём моки для демонстрации
        IpFilter mockFilter = mock(IpFilter.class);

        // Используем lenient() чтобы избежать UnnecessaryStubbingException
        lenient().when(mockFilter.contains(anyString())).thenReturn(false);
        lenient().when(mockFilter.getMemoryUsageBytes()).thenReturn(100L);
        lenient().when(mockFilter.getSize()).thenReturn(0);

        // Проверяем что моки работают
        boolean result = mockFilter.contains("1.1.1.1");
        assertFalse(result);

        // Verify вызов
        verify(mockFilter, atLeastOnce()).contains(anyString());
    }

    /**
     * Mock Type 2: Verify interactions (verify)
     */
    @Test
    void testMockVerify_Interactions() {
        IpFilter mockFilter = mock(IpFilter.class);

        when(mockFilter.contains("2.2.2.2")).thenReturn(false);
        when(mockFilter.contains("3.3.3.3")).thenReturn(true);

        // Вызываем методы
        mockFilter.contains("2.2.2.2");
        mockFilter.contains("3.3.3.3");

        // Проверяем вызовы
        verify(mockFilter, times(1)).contains("2.2.2.2");
        verify(mockFilter, times(1)).contains("3.3.3.3");
        verify(mockFilter, never()).clear();
    }

    /**
     * Mock Type 3: Spy (частичный мок)
     */
    @Test
    void testMockSpy_Usage() {
        // Создаём spy списка
        List<String> spyList = spy(new ArrayList<>(List.of("A", "B")));

        spyList.add("C");

        // Проверяем что метод был вызван
        verify(spyList).add("C");
        assertEquals(3, spyList.size());

        // Проверяем оригинальное поведение
        assertEquals("A", spyList.get(0));
        assertEquals("B", spyList.get(1));
        assertEquals("C", spyList.get(2));
    }

    /**
     * Mock Type 4: Exception testing with mocks
     */
    @Test
    void testMockException_Throwing() throws IOException {
        JsonHandler mockHandler = mock(JsonHandler.class);

        // Настраиваем исключение
        doThrow(new IOException("File not found"))
                .when(mockHandler).loadInput(anyString());

        // Проверяем что исключение выбрасывается
        assertThrows(IOException.class, () -> {
            mockHandler.loadInput("nonexistent.json");
        });
    }

    // ========================================================================
    // HAMCREST MATCHERS
    // ========================================================================

    /**
     * Минимум 2 различных вида Hamcrest matchers
     */
    @Test
    void testHamcrestMatchers() {
        OutputData result = service.processAndCompare(List.of("1.1.1.1"));
        FilterStatistics stats = result.getStatistics();

        // Matcher 1: greaterThan
        assertThat(stats.getHashSetMemoryBytes(), greaterThanOrEqualTo(0L));

        // Matcher 2: lessThan
        assertThat(stats.getFalsePositiveProbability(), lessThanOrEqualTo(1.0));

        // Matcher 3: equalTo
        assertThat(result.getResults().size(), equalTo(1));

        // Matcher 4: notNullValue
        assertThat(result.getTimestamp(), notNullValue());

        // Matcher 5: containsString (для recommendation)
        assertThat(result.getRecommendation(), not(emptyString()));
    }

    // ========================================================================
    // ADDITIONAL TESTS FOR MUTATION SCORE
    // ========================================================================

    @Test
    void testTimestamp_Format() {
        OutputData result = service.processAndCompare(List.of("1.1.1.1"));

        String timestamp = result.getTimestamp();
        assertNotNull(timestamp);
        // Проверяем формат ISO (должен содержать дату и время)
        assertTrue(timestamp.length() >= 19, "Timestamp должен быть в формате ISO");
        assertTrue(timestamp.contains("-") || timestamp.contains("T"),
                "Timestamp должен содержать разделители даты/времени");
    }

    @Test
    void testResult_IPField() {
        String testIp = "192.168.100.1";
        OutputData result = service.processAndCompare(List.of(testIp));

        assertEquals(testIp, result.getResults().get(0).getIp());
    }

    @Test
    void testStatistics_FalsePositiveProbability() {
        OutputData result = service.processAndCompare(List.of("1.1.1.1"));
        FilterStatistics stats = result.getStatistics();

        assertTrue(stats.getFalsePositiveProbability() >= 0);
        assertTrue(stats.getFalsePositiveProbability() <= 1);
    }

    @Test
    void testStatistics_TimeNonNegative() {
        OutputData result = service.processAndCompare(List.of("1.1.1.1"));
        FilterStatistics stats = result.getStatistics();

        assertTrue(stats.getHashSetTimeMs() >= 0);
        assertTrue(stats.getBloomFilterTimeMs() >= 0);
    }

    @Test
    void testStatistics_MemoryComparison() {
        OutputData result = service.processAndCompare(List.of("1.1.1.1"));
        FilterStatistics stats = result.getStatistics();

        // Оба значения памяти должны быть положительными
        assertTrue(stats.getHashSetMemoryBytes() >= 0);
        assertTrue(stats.getBloomFilterMemoryBytes() > 0);
    }

    @Test
    void testServiceConstructor_Default() {
        FilterService newService = new FilterService();
        assertNotNull(newService);
    }

    @Test
    void testProcessAndCompare_SingleIp() {
        OutputData result = service.processAndCompare(List.of("8.8.8.8"));

        assertEquals(1, result.getResults().size());
        assertNotNull(result.getStatistics());
        assertNotNull(result.getRecommendation());
        assertFalse(result.getRecommendation().isEmpty());
    }

    @Test
    void testProcessAndCompare_MultipleIps() {
        List<String> ips = List.of("1.1.1.1", "2.2.2.2", "3.3.3.3");
        OutputData result = service.processAndCompare(ips);

        assertEquals(3, result.getResults().size());

        // Проверяем каждый результат
        for (FilterResult fr : result.getResults()) {
            assertNotNull(fr.getIp());
            assertNotNull(fr.isMatch());
        }
    }

    @Test
    void testProcessAndCompare_StatisticsPopulated() {
        OutputData result = service.processAndCompare(List.of("1.1.1.1"));
        FilterStatistics stats = result.getStatistics();

        assertTrue(stats.getHashSetMemoryBytes() >= 0);
        assertTrue(stats.getBloomFilterMemoryBytes() >= 0);
        assertTrue(stats.getHashSetTimeMs() >= 0);
        assertTrue(stats.getBloomFilterTimeMs() >= 0);
        assertTrue(stats.getFalsePositiveProbability() >= 0);
    }

    @Test
    void testProcessAndCompare_ResultMatch() {
        OutputData result = service.processAndCompare(List.of("1.1.1.1"));
        FilterResult filterResult = result.getResults().get(0);

        // Для IP не в blacklist оба должны вернуть false
        assertFalse(filterResult.isHashSetResult());
        assertFalse(filterResult.isBloomFilterResult());
        // match должен быть true (оба false)
        assertTrue(filterResult.isMatch());
    }

    @Test
    void testOutputData_AllFieldsNotNull() {
        OutputData result = service.processAndCompare(List.of("1.1.1.1"));

        assertAll("All fields not null",
                () -> assertNotNull(result.getTimestamp()),
                () -> assertNotNull(result.getResults()),
                () -> assertNotNull(result.getStatistics()),
                () -> assertNotNull(result.getRecommendation())
        );
    }

    @Test
    void testFilterResult_AllFields() {
        OutputData result = service.processAndCompare(List.of("8.8.8.8"));
        FilterResult filterResult = result.getResults().get(0);

        assertAll("FilterResult fields",
                () -> assertNotNull(filterResult.getIp()),
                () -> assertTrue(filterResult.isHashSetResult() || !filterResult.isHashSetResult()),
                () -> assertTrue(filterResult.isBloomFilterResult() || !filterResult.isBloomFilterResult()),
                () -> assertTrue(filterResult.isMatch() || !filterResult.isMatch())
        );
    }

    @Test
    void testResult_MatchField() {
        OutputData result = service.processAndCompare(List.of("1.1.1.1"));
        FilterResult filterResult = result.getResults().get(0);

        // match должен быть true когда результаты совпадают
        boolean expectedMatch = (filterResult.isHashSetResult() == filterResult.isBloomFilterResult());
        assertEquals(expectedMatch, filterResult.isMatch());
    }

    @Test
    void testProcessAndCompare_DuplicateIPs() {
        // Тестируем дубликаты IP
        OutputData result = service.processAndCompare(
                List.of("1.1.1.1", "1.1.1.1", "1.1.1.1")
        );

        // Должно обработать все 3 IP (даже дубликаты)
        assertEquals(3, result.getResults().size());
    }

    @Test
    void testProcessAndCompare_EmptyString() {
        // Тестируем обработку пустой строки
        OutputData result = service.processAndCompare(List.of(""));

        assertNotNull(result);
        assertEquals(1, result.getResults().size());
    }

    @Test
    void testProcessAndCompare_NullIP() {
        // Тестируем обработку null IP
        List<String> ips = new ArrayList<>();
        ips.add(null);

        assertDoesNotThrow(() -> {
            OutputData result = service.processAndCompare(ips);
            assertNotNull(result);
            assertEquals(1, result.getResults().size());
        });
    }

    @Test
    void testRecommendation_MemoryDifference() {
        OutputData result = service.processAndCompare(
                List.of("1.1.1.1", "2.2.2.2", "3.3.3.3", "4.4.4.4", "5.5.5.5")
        );

        assertNotNull(result.getRecommendation());
        assertTrue(result.getRecommendation().length() > 0);
    }

    @Test
    void testRecommendation_LargeDataSet() {
        List<String> ips = new ArrayList<>();
        for (int i = 0; i < 50; i++) {
            ips.add(i + "." + i + "." + i + "." + i);
        }

        OutputData result = service.processAndCompare(ips);

        assertNotNull(result.getRecommendation());
        assertFalse(result.getRecommendation().isEmpty());
    }

    @Test
    void testProcessAndCompare_MixedResults() {
        OutputData result = service.processAndCompare(
                List.of("1.1.1.1", "2.2.2.2", "3.3.3.3")
        );

        assertEquals(3, result.getResults().size());

        // Проверяем что все результаты имеют IP
        for (FilterResult fr : result.getResults()) {
            assertNotNull(fr.getIp());
        }
    }

    @Test
    void testProcessAndCompare_InvalidIP() {
        assertDoesNotThrow(() -> {
            OutputData result = service.processAndCompare(List.of("invalid.ip.format"));
            assertNotNull(result);
            assertEquals(1, result.getResults().size());
        });
    }

    @Test
    void testServiceConstructor_WithMocks() {
        // Тест конструктора с моками
        IpFilter mockHashSet = mock(IpFilter.class);
        IpFilter mockBloom = mock(IpFilter.class);
        JsonHandler mockJsonHandler = mock(JsonHandler.class);

        FilterService serviceWithMocks = new FilterService(mockHashSet, mockBloom, mockJsonHandler);
        assertNotNull(serviceWithMocks);
    }
}
\end{verbatim}